\section{Síntesis y ampliación}

Esta clase parte de la forma abstracta de la AGI y termina en un mapa de sistemas que se puede ejecutar. Lo esencial de un Agent no es que «el modelo sepa invocar herramientas», sino que objetivos, estado, acciones, permisos, evaluación, aprendizaje, memoria, comunicación y recuperación formen un ciclo cerrado. Si falta cualquiera de estas capas, una demostración de capacidades se convierte en una automatización imposible de mantener.

Además, esta clase separa con claridad la «inteligencia del modelo» de la «confiabilidad del sistema»: la primera determina el límite superior de las acciones candidatas; la segunda determina qué acciones pueden llegar al mundo real. Un modelo más potente puede reducir algunos errores de razonamiento, pero no aporta por sí mismo auditoría, autorización, reversión, control de versiones ni responsables. Al estudiar sistemas de Agent, toda mejora del modelo debe evaluarse siempre dentro del plano de control completo.

\subsection{Una tabla que articula todo el Agent stack}

\begin{center}
\begin{tabularx}{\textwidth}{>{\bfseries}p{0.16\textwidth} X X X}
Capa & Pregunta central & Mecanismos representativos & Señales de falla \\
\hline
Goal & Qué quiere realmente el usuario & Análisis de restricciones, confirmación & Deriva del objetivo, pérdida de restricciones implícitas \\
Policy & Qué hacer en el siguiente paso & planning, MCTS, routing & Comportamiento voraz, ciclos, exceso de permisos \\
Environment & Dónde se ejecutan las acciones & browser/API/tools & Cambios en la página, acciones irreversibles \\
Evaluation & Cómo saber que se hizo bien & REAL, slice, trace checks & La puntuación global oculta la cola \\
Learning & Cómo mejorar a partir de las trayectorias & critique, DPO, RL & Explotación de la recompensa, datos contaminados \\
Memory & Cómo conservar el estado entre pasos y sesiones & scratchpad, retrieval, profile & Caducidad, filtraciones, persistencia de errores \\
Coordination & Cómo colaboran varios Agent & hierarchy, sync, MCP/A2A & Conflictos, duplicación, esperas \\
Deployment & Cómo asumir responsabilidad en el mundo real & budget, audit, override & Pérdida de control, falta de trazabilidad, imposibilidad de recuperación \\
\end{tabularx}
\end{center}

\begin{importantbox}{La síntesis más importante de esta clase}
La forma de producto de la AGI aún no está definida, pero las restricciones de ingeniería de un Agent confiable ya son claras: \textbf{evaluar de forma reproducible, aprender sobre trayectorias, actuar según permisos, recordar con fuentes, colaborar mediante protocolos y, ante una falla, detenerse y devolver el control.}
\end{importantbox}

\subsection{Doce preguntas antes de iniciar un proyecto de Agent}

\begin{multicols}{2}
\small
\begin{enumerate}
  \item ¿Puede una máquina determinar el objetivo del usuario y la condición de término?
  \item ¿Qué acciones son reversibles y cuáles requieren confirmación?
  \item ¿El Agent actúa mediante una API o controla directamente la interfaz?
  \item ¿El benchmark cubre sitios web reales, la longitud de las tareas y los slice de riesgo?
  \item ¿Se registra la trace completa de observation--action--result?
  \item ¿La taxonomy de fallas permite ubicar el origen en el modelo, las herramientas, la memory o la policy?
  \item ¿La búsqueda y la self-critique tienen costo y condición de paro?
  \item ¿Los preference data incluyen tanto ramas exitosas como fallidas?
  \item ¿La memoria de largo plazo tiene fuente, caducidad, permisos y edición por parte del usuario?
  \item ¿El sistema multi agent realmente puede paralelizarse o solo agrega comunicación?
  \item Más allá de MCP/A2A, ¿cómo se gestionan la identidad, la autorización y las versiones?
  \item Tras el lanzamiento, ¿quién monitorea, quién toma el control y cómo se revierte?
\end{enumerate}
\normalsize
\end{multicols}

\subsection{Lecturas adicionales}

\begin{itemize}
  \item Putta et al., \emph{Agent Q: Advanced Reasoning and Learning for Autonomous AI Agents}: MCTS, self-critique y aprendizaje por preferencias. \href{https://arxiv.org/abs/2408.07199}{arXiv:2408.07199}
  \item AGI, Inc., \emph{REAL: Benchmarking Autonomous Agents on Deterministic Simulations of Real Websites}: entornos web deterministas y leaderboard. \href{https://github.com/agi-inc/real}{REAL repository}
  \item Lilian Weng, \emph{LLM Powered Autonomous Agents}: panorama de la arquitectura de memory, planning, tools y action. \href{https://lilianweng.github.io/posts/2023-06-23-agent/}{article}
  \item Documentación oficial de Model Context Protocol: protocolo abierto de conexión para herramientas, recursos y prompt. \href{https://modelcontextprotocol.io/}{official docs}
\end{itemize}

\paragraph{Alcance de la evidencia.}
Esta clase registra el contenido impartido en abril de 2025. REAL, AgentQ, MultiOn, MCP y los resultados de los modelos de ese momento deben entenderse según la fecha de la clase; estas notas conservan sus mecanismos y sus lecciones de ingeniería, sin convertir las demo, el estado de los productos ni las cifras de benchmark en garantías vigentes de productos en 2026.

\end{document}
